\clearpage\SourceLeaf{197}\ArithmeticHead{112}{L}{L. SCHLESINGER, ÜBER GAUSS’ FUNKTIONENTHEORETISCHE ARBEITEN.}
Es wurde bereits darauf hingewiesen (oben S. 90), dass \textsc{Gauss} sehr bald
über die Theorie der elliptischen Funktionen hinausgegangen ist, indem er
die allgemeineren Transzendenten, die durch die \textsc{Gausssche} Differentialgleichung
definiert werden, in den Kreis seiner Betrachtungen zog.

Während also bei \textsc{Abel} und \textsc{Jacobi} die Theorie der elliptischen \textls[100]{Integrale}
im Vordergrunde steht und demgemäss der weitere Fortschritt nach
der Seite der Integrale der allgemeinen algebraischen Funktionen hin erfolgt,
wird \textsc{Gauss} durch seinen in der Theorie des agM. wurzelnden Gedankengang
von der Theorie der elliptischen Funktionen nach der Seite der Theorie der
linearen Differentialgleichungen zweiter Ordnung hin gelenkt und erscheint
so als Bahnbrecher und Vorläufer für die funktionentheoretische Forschung
der letzten Dezennien des neunzehnten Jahrhunderts, wie sie durch die Namen
\textsc{Fuchs}, \textsc{Poincaré} und \textsc{Klein} charakterisiert wird.

Dass \textsc{Gauss} nicht im Besitz von Methoden war, um Integrale von mehrdeutigen
Funktionen im komplexen Gebiete zu untersuchen, hat \textsc{Riemann} sofort
durchschaut, als er (vergl. oben S. 5, Fussnote \textsuperscript{5)}) sich mit dem Nachlass
beschäftigt hatte. In der Tat sagt er in der Einleitung der 1857 erschienenen
Abhandlung \textit{Beiträge zur Theorie der durch die Gausssche Reihe darstellbaren
Funktionen}\footnote{\textsc{Riemanns} Werke, 1892, Abhandlung IV, S. 67.}: »Die \textsc{Gausssche} Reihe \(F(\alpha,\beta,\gamma,x)\), als Funktion ihres vierten
Elements \(x\) betrachtet, stellt diese Funktion nur dar, so lange der Modul von
\(x\) die Einheit nicht überschreitet. Um diese Funktion in ihrem ganzen Umfange,
bei unbeschränkter Veränderlichkeit dieses ihres Arguments zu untersuchen,
bieten die bisherigen Arbeiten über dieselbe zwei Wege dar. Man
kann nämlich entweder von einer linearen Differentialgleichung, welcher sie
genügt, ausgehn, oder von ihrem Ausdruck durch bestimmte Integrale. Jeder
dieser Wege gewährt eigentümliche Vorteile; jedoch ist bisher in der reichhaltigen
Abhandlung von \textsc{Kummer} im 15. Bande des mathematischen Journals
von \textsc{Crelle} und auch in den noch unveröffentlichten Untersuchungen von
\textsc{Gauss} nur der erste betreten, wohl hauptsächlich deshalb, weil die Rechnung
mit bestimmten Integralen zwischen komplexen Grenzen noch zu wenig ausgebildet
war, oder doch nicht als einem grossen Leserkreise geläufig vorausgesetzt
werden konnte.« — Diesen zweiten Weg zum Studium der gedachten\par
\clearpage\SourceLeaf{198}\ArithmeticHead{113}{R}{Vd. ANWENDUNGEN. ZUSAMMENFASSUNG.}
\noindent Funktion hat \textsc{Riemann} selbst beschritten in einer Vorlesung \textit{Über Funktionen
einer veränderlichen Grösse, insbesondere über hypergeometrische Reihen
und verwandte Transzendenten}, die er im Wintersemester 1858/59 gehalten hat
und von der einige Teile 1902 veröffentlicht worden sind\footnote{Siehe Nachträge zu \textsc{Riemanns} Werken, herausgegeben von M. \textsc{Noether} und W. \textsc{Wirtinger}, 1902, III, S. 69 ff. Vergl. auch Vorwort, S. v, wo über das dieser Veröffentlichung zugrunde liegende Heft berichtet wird.}. In dieser Veröffentlichung
ist für uns auch noch der letzte Artikel (a. a. O. S. 91) von Interesse,
wo \textsc{Riemann} von den vollständigen elliptischen Integralen \(K,K'\) (in \textsc{Jacobischer}
Bezeichnung) ausgehend, eine Theorie der Modulfunktion \(k^2\) als Funktion
von \(K'/K\) skizziert. Er gibt am Schluss (a. a. O. S. 93) die Figur des Fundamentalbereichs,
die, da ja \(K'/K\) mit dem \textsc{Gaussschen} \(t\) übereinstimmt, genau
die Form hat, wie bei \textsc{Gauss} (siehe oben Figur 2). Da \textsc{Riemann} den \textsc{Gaussschen}
Nachlass gekannt hat, ist es wohl möglich, dass er diese Figur bei
\textsc{Gauss} kennen gelernt hat; jedenfalls ist es bemerkenswert, dass sie noch zehn
Jahre vor ihrer Publikation in \textsc{Gauss’} Werken, durch \textsc{Riemann} auf die Tafel
eines Göttinger Hörsaals hingezeichnet worden ist. Es muss hier überhaupt
die Tatsache festgestellt werden, dass der \textsc{Gausssche} Nachlass zur Theorie der
elliptischen Funktionen verhältnismässig wenig Interesse gefunden hat. Man
kann wohl sagen, dass es fast ausschliesslich infolge der durch drei Göttinger
Professoren, \textsc{Schering}, \textsc{Fuchs} und \textsc{Klein}, gegebenen Anregungen dazu gekommen
ist, dass Studien über diesen Nachlass angestellt worden sind, wobei
noch in Betracht kommt, dass bis zum Tode \textsc{Scherings} (1897) diese Studien
nur auf den Band III der Werke, nicht aber auf die Handschriften selbst
gegründet werden konnten\footnote{Siehe die Bibliographie der auf den Nachlass bezüglichen Publikationen bei \textsc{Geppert} a. a. O. \textsuperscript{83)}, S. 190—194. Wir bemerken, dass die offenbar durch \textsc{Hermite} veranlasste, sehr verdienstliche Abhandlung des P. \textsc{Pepin}, \textit{Introduction à la théorie des fonctions elliptiques d’après les œuvres posthumes de Gauss}, Memorie della Pont. Accad. dei Nuovi Lincei 9, 1893, S. 1—129, wohl auch mittelbar auf eine Anregung \textsc{Scherings} zurückgeht, der 1877, als \textsc{Hermite} zur Jahrhundertfeier von \textsc{Gauss’} Geburtstag in Göttingen weilte, diesen auf die Schätze hingewiesen haben mag, die der in Werke III abgedruckte Nachlass birgt.}.

Inwieweit \textsc{Gauss} seine Untersuchungen zum Ausgangspunkt allgemeiner
funktionentheoretischer Betrachtungen genommen haben mag, lässt sich aus
seinem Nachlass nur unvollkommen feststellen, und das wenige an nachgelassenen
Notizen und Briefschaften, das uns einen Einblick gewährt in den\par
\HistoricalFoot{X2 Abh. 2.}{15}
\clearpage\SourceLeaf{199}\ArithmeticHead{114}{L}{L. SCHLESINGER, ÜBER GAUSS’ FUNKTIONENTHEORETISCHE ARBEITEN.}
\noindent Umfang seiner funktionentheoretischen Kenntnisse, stammt aus späterer Zeit,
so dass wir den Bericht über diesen Gegenstand auf die folgenden Abschnitte
verschieben müssen.
